% ============================================================================
\chapter{调度器}
\label{ch:scheduler}
% Stage D-9 ~ D-12
% ============================================================================

\section{本章目标}

\begin{itemize}
  \item 理解协作式与抢占式调度的区别
  \item 设计可插拔 \texttt{sched\_ops\_t} 接口，解耦调度算法
  \item 配置 PIT 8253 定时器产生周期中断（IRQ0）
  \item 实现上下文切换汇编（保存/恢复寄存器 + CR3 + TSS.ESP0）
  \item 实现 Round-Robin 调度器后端
  \item 实现优先级调度器后端（进阶）
  \item 通过 \texttt{KCONFIG\_SCHED\_BACKEND} 一键切换
\end{itemize}

% ============================================================================
\section{概念：CPU 调度}
\label{sec:sched-concept}
% ============================================================================

% TODO:
% - 调度的目标：多个就绪进程共享有限 CPU → 公平/高效/响应快
% - 协作式调度 (Cooperative)：进程自愿让出 CPU（yield）
%     优点：实现简单，无需定时器中断
%     缺点：恶意/有 bug 的进程可以霸占 CPU 永不释放
% - 抢占式调度 (Preemptive)：定时器中断强制切换
%     优点：内核保证公平性，用户程序无法垄断 CPU
%     缺点：需要定时器 + 更复杂的上下文切换
% - 我们实现抢占式调度（Unix/Linux 标准方式）
% - 时间片 (Time Quantum)：每个进程连续运行的最大时间
%     太短 → 切换开销大；太长 → 响应慢
%     典型值：10ms (100Hz) ~ 100ms (10Hz)
% - 图：三个进程的甘特图（时间片轮转示意）

\subsection{调度算法分类}

% TODO:
% - Round-Robin（轮转）：简单公平，所有进程等优先级
% - Priority Scheduling（优先级）：高优先级先运行，可能饥饿
% - MLFQ (Multi-Level Feedback Queue)：多级反馈，Linux CFS 的前身
% - CFS (Completely Fair Scheduler)：Linux 2.6.23+，基于 vruntime 红黑树
% - 我们先实现 RR，再进阶到优先级调度

% ============================================================================
\section{可插拔接口：\texttt{sched\_ops\_t}}
\label{sec:sched-interface}
% ============================================================================

% TODO:
% - 接口设计（与 pmm_ops_t 同一模式）：
%     typedef struct sched_ops {
%         const char* name;                           // "round_robin", "priority"
%         void (*init)(void);                          // 初始化就绪队列
%         void (*enqueue)(struct task_struct* task);   // 加入就绪队列
%         void (*dequeue)(struct task_struct* task);   // 从就绪队列移除
%         struct task_struct* (*pick_next)(void);      // 选择下一个运行的进程
%         void (*tick)(struct task_struct* current);   // 定时器 tick 回调
%     } sched_ops_t;
% - dispatch 层：sched.c 持有 g_sched_ops 全局指针
%     sched_register_backend() → 注册
%     sched_init() → 初始化
%     schedule() → 核心调度函数：
%         1. 当前进程入队（如果还是 READY）
%         2. pick_next() 选出下一个
%         3. context_switch(current, next)
% - kconfig.h: #define KCONFIG_SCHED_BACKEND 0  // 0=round_robin, 1=priority

% ============================================================================
\section{PIT 定时器配置}
\label{sec:sched-pit}
% ============================================================================

% TODO:
% - PIT 8253/8254 硬件：3 个独立计数器，我们用 Counter 0 (IRQ0)
% - 基础频率：1193182 Hz
% - 配置目标频率 100Hz（每 10ms 一次 tick）：
%     divisor = 1193182 / 100 = 11932
%     outb(0x43, 0x36)  // Counter 0, lo/hi byte, mode 3 (square wave)
%     outb(0x40, divisor & 0xFF)   // low byte
%     outb(0x40, divisor >> 8)     // high byte
% - [WHY] 100Hz 是经典选择：10ms tick 兼顾响应性和切换开销
% - IRQ0 handler 调用 sched_tick() → 可能触发 schedule()
% - pic_clear_mask(0) 启用 IRQ0

% ============================================================================
\section{上下文切换}
\label{sec:sched-context-switch}
% ============================================================================

% TODO:
% - 上下文切换 = 保存当前进程状态 + 恢复下一个进程状态
% - context_switch(prev, next) 汇编实现：
%     ; === 保存 prev ===
%     pushad
%     push ds / es / fs / gs
%     mov [prev->esp], esp          ; 保存内核栈指针
%     ; === 切换地址空间 ===
%     mov eax, [next->page_directory]
%     mov cr3, eax                   ; TLB 全刷新
%     ; === 更新 TSS ===
%     mov eax, [next->kernel_stack]
%     mov [tss.esp0], eax            ; 下次中断用 next 的内核栈
%     ; === 恢复 next ===
%     mov esp, [next->esp]           ; 切换到 next 的内核栈
%     pop gs / fs / es / ds
%     popad
%     ret                            ; 返回到 next 上次被切走的位置

\begin{cpustate}
上下文切换后：
\begin{itemize}
  \item CR3 指向新进程的页目录 → 地址空间完全切换
  \item TSS.ESP0 指向新进程的内核栈 → 下次中断用正确的栈
  \item ESP 指向新进程上次保存的栈帧 → \texttt{popad} 恢复所有通用寄存器
  \item \texttt{ret} 跳到新进程上次被 \texttt{schedule()} 切走时的下一条指令
\end{itemize}
\end{cpustate}

% ============================================================================
\section{后端一：Round-Robin 调度器}
\label{sec:sched-rr}
% ============================================================================

% TODO:
% - 数据结构：环形双向链表（就绪队列）
%     head → task_A ↔ task_B ↔ task_C → head
% - enqueue：插入队尾（O(1)）
% - dequeue：从链表移除（O(1)）
% - pick_next：取队头（O(1)）
% - tick：current->remaining_ticks--
%     如果 == 0 → 设 need_resched = 1 → 重置 remaining_ticks = TIME_SLICE
% - TIME_SLICE = 10（10 * 10ms = 100ms 时间片）

% ============================================================================
\section{后端二：优先级调度器（进阶）}
\label{sec:sched-priority}
% ============================================================================

% TODO:（进阶，可跳过）
% - 数据结构：多级就绪队列 ready_queue[MAX_PRIORITY]
%     优先级 0（最高）~ N（最低）
% - pick_next：从最高优先级非空队列取队头
% - 防饥饿：aging 机制
%     每次 tick，所有等待中的进程优先级提升 1 级
% - KCONFIG_SCHED_BACKEND=1 切换

% ============================================================================
\section{验证}
\label{sec:sched-verify}
% ============================================================================

% TODO:
% - 测试 1：两个用户进程各打印 A/B → 串口看到交替输出 ABABAB...
% - 测试 2：三个进程不同工作量 → 观察公平分配
% - 测试 3：进程 exit → 自动从就绪队列移除 → 不再被调度
% - 测试 4（优先级）：高优先级抢占低优先级
% - 里程碑 D 集成测试：fork+exec+scheduler 完整流程
% - shell 命令：sched info / sched test

% ============================================================================
\section{本章小结}
% ============================================================================

调度器通过可插拔 \texttt{sched\_ops\_t} 接口，将调度算法与内核核心解耦。
Round-Robin 提供了最简单的公平调度，优先级调度允许差异化的 CPU 分配。
配合 PIT 定时器的抢占式设计，保证了任何用户进程都无法独占 CPU。
至此"进程与用户态"里程碑 D 全部完成。下一部分进入文件系统。
